% ============================================================
%  MEAN-RECOVERY AS THE UNIVERSAL PHYSICAL LAW:
%  FROM PARTITION STATES TO ELECTROMAGNETIC SIGNALS,
%  WITH THE UNIFIED SIGNAL HYPOTHESIS AS A THEOREM
% ============================================================

\documentclass[twocolumn,10pt,a4paper]{article}

\usepackage[utf8]{inputenc}
\usepackage[T1]{fontenc}
\usepackage{lmodern}
\usepackage[a4paper,top=2cm,bottom=2.2cm,left=1.8cm,right=1.8cm,columnsep=0.6cm]{geometry}
\usepackage{amsmath,amssymb,amsthm,mathtools}
\usepackage{bm}
\usepackage{booktabs}
\usepackage{array}
\usepackage{enumitem}
\usepackage{microtype}
\usepackage[round,sort&compress]{natbib}
\usepackage[colorlinks=true,linkcolor=blue!60!black,
            citecolor=blue!60!black,urlcolor=blue!60!black]{hyperref}
\usepackage{fancyhdr}

\pagestyle{fancy}
\fancyhf{}
\fancyhead[LE,RO]{\thepage}
\fancyhead[RE]{\small Mean-Recovery as the Universal Physical Law}
\fancyhead[LO]{\small Sachikonye}
\renewcommand{\headrulewidth}{0.4pt}

\newtheoremstyle{thmstyle}{\topsep}{\topsep}{\itshape}{}{\bfseries}{.}{.5em}{}
\theoremstyle{thmstyle}
\newtheorem{theorem}{Theorem}[section]
\newtheorem{lemma}[theorem]{Lemma}
\newtheorem{proposition}[theorem]{Proposition}
\newtheorem{corollary}[theorem]{Corollary}
\theoremstyle{definition}
\newtheorem{definition}[theorem]{Definition}
\newtheorem{axiom}{Axiom}
\newtheorem{example}[theorem]{Example}
\theoremstyle{remark}
\newtheorem{remark}[theorem]{Remark}

\newcommand{\R}{\mathbb{R}}
\newcommand{\N}{\mathbb{N}}
\newcommand{\kB}{k_{\mathrm{B}}}
\newcommand{\nP}{n_{\mathrm{P}}}
\newcommand{\Parts}{\mathcal{P}}
\newcommand{\Sv}{\mathbf{S}}
\newcommand{\Sk}{S_{k}}
\newcommand{\St}{S_{t}}
\newcommand{\Se}{S_{e}}
\newcommand{\Hologram}{H}
\newcommand{\Signal}{\mathbf{S}}

% =========================================================
\title{\textbf{Mean-Recovery as the Universal Physical Law:\\
From Partition States to Electromagnetic Signals,\\
with the Unified Signal Hypothesis as a Theorem}}

\author{
Kundai Farai Sachikonye\\
Technical University of Munich\\
Chair of Proteomics and Bioanalytics\\
\texttt{kundai.sachikonye@wzw.tum.de}
}
\date{\today}
% =========================================================

\begin{document}
\maketitle

% =========================================================
\begin{abstract}
\noindent
The \emph{Unified Signal Hypothesis} \citep{companion:signal_source} conjectures
that a single physical electromagnetic beam simultaneously encodes all
frequency-domain decompositions (radio through optical), with each providing
an equally valid interpretation constrained by a universal mean-recovery
principle.  This paper proves the conjecture as a theorem within the
partition algebra of bounded phase space.

The partition algebra provides what the $L^2$ Hilbert space framework
cannot: a \emph{discrete}, \emph{countable}, \emph{physically grounded}
version of the same structure.  We prove four results.

(i) \textbf{Unified Signal Theorem}
(Theorem~\ref{thm:unified_signal}): any physical signal generated by a
bounded oscillatory system simultaneously admits all valid frequency-band
decompositions, with mean-recovery as the universal constraint.  This
follows from the Bounded Phase Space Axiom + Poincar\'e recurrence: the
mean-recovery constraint is the formal statement that every trajectory
returns to its origin, so all decompositions must share the same mean
(the physical state).

(ii) \textbf{Frequency as Derived Quantity Theorem}
(Theorem~\ref{thm:frequency_derived}): all electromagnetic frequencies are
derived from the single counting identity $dM/dt = \omega/(2\pi) = 1/\langle\tau_p\rangle$.
There is one fundamental rate (partition state traversal per unit time);
all frequencies are its multi-dimensional projections.

(iii) \textbf{Cross-Domain Coherence from Selection Rules}
(Theorem~\ref{thm:coherence}): the exponential decay of inconsistency between
different frequency-band decompositions, $C(D_1,D_2) \sim e^{-\delta/B}$,
is the partition-algebra statement that transitions violating selection rules
are structurally forbidden (type errors), not merely rare.  Spectrally
distant bands agree because accessing both requires a selection-rule
violation.

(iv) \textbf{Four Synonyms, One Structure}
(Theorem~\ref{thm:synonyms}): geometric aperture (optics), virtual substate
(partition algebra), chemical transition state (chemistry), off-shell
propagator (QFT), and off-band virtual component (signal theory) are five
names for the same mathematical object --- an intermediate decomposition
that lies outside the physical domain but whose mean satisfies the physical
constraint.

We trace the mean-recovery constraint through five physical domains
(MS fragmentation, LHC collisions, molecular spectral holography,
atomic spectroscopy, and single-beam positioning) and show that in each
case the ``constraint'' is not an imposed law but the direct consequence
of bounded phase space.

\medskip
\noindent\textbf{Keywords:} mean-recovery; unified signal hypothesis;
bounded phase space; frequency as derived quantity; cross-domain coherence;
virtual substates; partition algebra; five synonyms; spectral hologram;
signal unification
\end{abstract}

\tableofcontents

% =========================================================
\section{Introduction}
\label{sec:intro}
% =========================================================

\subsection{The conjecture and its gap}

A recent paper (the Unified Signal Hypothesis, \citealt{companion:signal_source})
proposes that all frequency-domain decompositions of a single electromagnetic
signal are simultaneously valid, constrained by a universal mean-recovery
law.  The key results are established for $L^2(\R)$ functions:
existence of off-shell decompositions (Miracle Principle); universal
mean-recovery linking all domain-specific interpretations; continuity of
domain transformations; multi-domain positioning.

The framework leaves four things as conjectures or ungrounded claims:
\begin{enumerate}[nosep]
\item Why does mean-recovery hold universally?  The $L^2$ proof shows
  existence but not necessity.
\item Why does frequency feel intrinsic?  The paper offers intuition but
  no proof.
\item Why do distant frequency bands agree exponentially well?  The
  cross-domain coherence bound $C \sim e^{-\delta/B}$ is asserted without
  derivation.
\item What is the fundamental constant of physics?  The paper conjectures
  it is the mean-recovery constraint, not $c$.
\end{enumerate}

The partition algebra answers all four precisely because it is the discrete,
countable, physically grounded version of the same structure.

\subsection{The partition algebra as the discrete signal theory}

The partition bijection $\PhiMap: \mathbb{Z}^+ \to \Parts$ is the
discrete analog of the Fourier transform: it maps every positive integer
(every oscillation count) to a unique partition state (a unique frequency
component).  The S-entropy embedding maps partition states to $[0,1]^3$
exactly as the Fourier transform maps time-domain signals to frequency-domain
amplitudes.  Virtual substate decompositions with mean-recovery are the
discrete analog of off-shell frequency decompositions.

The $L^2$ theory is the \emph{continuum limit} of the partition algebra.
The partition algebra answers the four conjectures because it operates on
the discrete foundation from which the continuum theory emerges.

% =========================================================
\section{The Unified Signal Theorem}
\label{sec:unified}
% =========================================================

\begin{axiom}[Bounded Phase Space]
Every persistent physical system occupies a bounded region of phase space
with finite Liouville measure.
\end{axiom}

\begin{theorem}[Poincar\'e Return Implies Mean-Recovery]
\label{thm:poincare_mean_recovery}
For any bounded oscillatory system and any virtual substate decomposition
$(\Sv_1,\ldots,\Sv_N)$ of the system's partition state $\Sv$, the
mean-recovery constraint $N^{-1}\sum_i \Sv_i = \Sv$ is not an imposed law
but the direct consequence of Poincar\'e recurrence: every trajectory must
return to the physical state $\Sv$.
\end{theorem}

\begin{proof}
By the Bounded Phase Space Axiom + Poincar\'e recurrence
\citep{poincare1890}, every trajectory $\Sv(t)$ returns arbitrarily close
to its initial value $\Sv(0)$.  A virtual substate decomposition
$(\Sv_1,\ldots,\Sv_N)$ parametrises $N$ simultaneous virtual "probes" of the
same oscillatory trajectory at $N$ different phases.  The time-average of
any bounded trajectory over a full oscillation period equals its mean value:
$\langle\Sv\rangle = N^{-1}\sum_i \Sv_i$.
Since Poincar\'e recurrence guarantees the trajectory returns to $\Sv$,
the mean must equal $\Sv$.  This is mean-recovery.
\end{proof}

\begin{theorem}[Unified Signal Theorem]
\label{thm:unified_signal}
A single physical signal generated by a bounded oscillatory system
simultaneously admits all valid frequency-band decompositions.
Each decomposition is equally valid; the physical state is the unique
invariant across all decompositions; and mean-recovery is the necessary
and sufficient constraint linking them.
\end{theorem}

\begin{proof}
\emph{Existence}: By the Miracle Principle \citep{companion:trajectory},
for any physical state $\Sv \in [0,1]^3$ and any $N-1$ arbitrary virtual
components, the unique $N$-th component satisfying mean-recovery exists.
Hence infinitely many decompositions exist.

\emph{Uniqueness of invariant}: The physical state $\Sv$ is the mean of
all components in any valid decomposition.  It is therefore invariant under
all domain transformations (which are changes of decomposition basis).

\emph{Necessity of mean-recovery}: By Theorem~\ref{thm:poincare_mean_recovery},
mean-recovery is the Poincar\'e recurrence condition in decomposition language.
Any decomposition that violates mean-recovery would correspond to a trajectory
that never returns to its initial value --- impossible in bounded phase space.
\end{proof}

\begin{corollary}[The Physical State is the Invariant, Not the Frequency]
The frequency of a signal is a property of the decomposition basis, not
of the physical state.  The physical state (partition address $\Sv \in [0,1]^3$)
is the unique invariant across all decompositions.
\end{corollary}

% =========================================================
\section{Frequency as a Derived Quantity}
\label{sec:frequency}
% =========================================================

\begin{theorem}[All Frequencies Are Counting]
\label{thm:frequency_derived}
Every electromagnetic frequency is a projection of the single counting
identity $dM/dt = \omega/(2\pi) = 1/\langle\tau_p\rangle$ onto a
measurement basis.  Frequency is not a primitive property of signals;
it is the rate of partition-state traversal projected onto a chosen
decomposition axis.
\end{theorem}

\begin{proof}
From the Time-Count Identity (Theorem~3.2 of \citealt{companion:math}):
the partition count $M(t)$ satisfies $dM/dt = \omega/(2\pi) = 1/\langle\tau_p\rangle$
for any bounded oscillator.  This is one equation, not a separate equation
for each frequency.

A measurement in frequency band $D = (\omega_c, \Delta\omega)$ extracts
the component of $dM/dt$ near $\omega_c$: it restricts the partition
traversal to states with oscillation frequency near $\omega_c$.
Different frequency bands are different projections of the same $dM/dt$.

The categorical angular units \citep{companion:math} make this explicit:
$c = 2\pi\,\text{rad/tick}$, $\hbar = E_\text{tick}/(2\pi)$,
$k_B = E_\text{tick}/\ln(d+1)$.
All three encode the same tick rate; the constants $c$, $\hbar$, $k_B$
are three projections of the single counting rate onto three measurement
bases (spatial, quantum, thermal).  Frequency, wavelength, and phase are
derived quantities in exactly the same way.
\end{proof}

\begin{remark}[Why frequency feels intrinsic]
Most real signals are narrowband: their partition-traversal rate is
concentrated near one frequency.  The decomposition appears unique because
it is so tightly constrained (Theorem~\ref{thm:coherence} below shows the
coherence bound forces agreement for spectrally distant bands).
Broadband signals (ultrashort pulses) reveal the decomposition ambiguity
directly: their frequency content is explicitly dependent on the measurement
window and decomposition choice.
\end{remark}

% =========================================================
\section{Cross-Domain Coherence from Selection Rules}
\label{sec:coherence}
% =========================================================

\begin{theorem}[Cross-Domain Coherence from Selection Rules]
\label{thm:coherence}
The exponential decay of inconsistency between frequency-band decompositions,
$C(D_1, D_2) \sim e^{-\delta/B}$ where $\delta$ is spectral separation
\citep{companion:signal_source}, is the partition-algebra statement that
transitions between partition states at spectral separation $\delta$
violate selection rules and are structurally forbidden.
\end{theorem}

\begin{proof}
A transition between partition states corresponding to frequency bands $D_1$
and $D_2$ requires $|\Delta n| = |\omega_{c,1}/\omega_0 - \omega_{c,2}/\omega_0|$
shell levels (where $\omega_0$ is the reference oscillator frequency).
In atomic spectroscopy, the electric-dipole selection rule enforces
$|\Delta\ell| = 1$: forbidden transitions are type errors in the vaHera
refinement type system \citep{companion:types}.

For a molecular or electromagnetic oscillator, the selection rule is
energetic rather than dipole-symmetry: transitions spanning $|\Delta n|$
shells require energy $\sim |\Delta n| \cdot E_\text{tick}$.
The probability of such a transition occurring in a given coherence time
is suppressed by the Boltzmann factor $e^{-|\Delta n| E_\text{tick}/k_B T}
= e^{-\delta/B}$ where $\delta \propto |\Delta n|$ is the spectral
separation and $B \propto k_B T / E_\text{tick}$ is the bandwidth constant.

Hence the cross-domain coherence bound is not an empirical observation
but the partition-algebra statement that spectrally distant decompositions
must agree because connecting them requires forbidden transitions.
\end{proof}

\begin{corollary}[Radio and Optical Bands Always Agree]
For radio ($\omega_c \sim 10^9$\,Hz) and optical ($\omega_c \sim 10^{15}$\,Hz),
$\delta \sim 6 \times 10^{15}$\,Hz.  With $B \sim 10^9$\,Hz:
$C \sim e^{-6\times 10^6} \approx 0$.  The two bands are forced to agree
to essentially infinite precision by the partition selection rules.
This is why all four LHC optical channels (spanning transition radiation
to synchrotron: 10\,keV to MeV range) give consistent particle identification:
their spectral separation is large, coherence bound is tiny.
\end{corollary}

% =========================================================
\section{Five Synonyms, One Structure}
\label{sec:synonyms}
% =========================================================

The Aperture Identification Theorem of the trajectory-completion paper
\citep{companion:trajectory} established five synonyms for the same object.
The unified signal framework adds a sixth.

\begin{theorem}[Six Synonyms, One Object]
\label{thm:synonyms}
The following six descriptions name the same mathematical structure:
an intermediate decomposition that lies outside the physical domain but
whose mean satisfies the physical constraint.

\begin{enumerate}[nosep,label=(\arabic*)]
\item \textbf{Geometric aperture} (optics): the region of allowed
  sub-coordinate decompositions of a fixed global coordinate, intersected
  with the off-shell region.
\item \textbf{Virtual substate} (partition algebra): an off-shell component
  $\Sv_i \notin [0,1]^3$ in a decomposition with mean $\Sv \in [0,1]^3$.
\item \textbf{Chemical transition state} (chemistry): an unstable intermediate
  configuration mediating bond cleavage, accessible only transiently
  \citep{companion:graph}.
\item \textbf{Off-shell Feynman propagator} (QFT): an internal line of a
  Feynman diagram with $q^2 \neq m^2c^4$ \citep{companion:virtual}.
\item \textbf{Miracle subtask} (information theory): a subtask with local
  S-value 100 (maximally misaligned) that contributes to a globally aligned
  composition \citep{companion:trajectory}.
\item \textbf{Off-band virtual component} (signal theory): a frequency-domain
  component outside the nominal band that contributes to the physical
  signal mean.
\end{enumerate}

All six share four properties: (i) they enable a transition;
(ii) they are not consumed; (iii) they transfer zero information about
themselves to the transition outcome; (iv) they are necessary for
sub-exponential complexity.
\end{theorem}

\begin{proof}
Properties (i)-(iv) for each synonym:
\begin{itemize}[nosep]
\item (1)--(2): identical by definition, proved in \citealt{companion:trajectory}.
\item (2)--(3): Theorem~\ref{thm:virtual_ts} of \citealt{companion:graph}.
\item (2)--(4): off-shell mass $p^2 - m^2$ = partition malformation energy
  at the virtual component, proved in \citealt{companion:virtual}.
\item (2)--(5): miracle subtask = decomposition with extreme local S-value,
  proved in \citealt{companion:trajectory}.
\item (2)--(6): the off-band component is outside $D$ (off-shell in frequency)
  but the mean is within the physical signal space.  Properties (i)-(iv)
  follow from Theorem~\ref{thm:unified_signal}: the component enables the
  decomposition (i), is not consumed by mean-recovery (ii), the physical
  state is the mean not the component itself (iii), and removing off-band
  components collapses to $\Theta(V)$ complexity (iv).
\end{itemize}
\end{proof}

% =========================================================
\section{Mean-Recovery in Five Physical Domains}
\label{sec:domains}
% =========================================================

We trace the mean-recovery constraint through five domains, showing that
in each case it is the Poincar\'e recurrence condition in disguise.

\subsection{MS/MS fragmentation}

The precursor ion oscillates in the Orbitrap trap.  Its fragment ions
appear as subharmonic frequency components $\omega_f = \omega_p\sqrt{m_p/m_f}$
in the precursor transient \citep{companion:tensor}.  This is mean-recovery:
the fragment and precursor are two decompositions of the same oscillatory
signal; their frequency ratio ensures the means are consistent.

\subsection{LHC collision events}

The $pp$ collision event produces signals in four detector subsystems
plus four optical channels \citep{companion:extensor}.  The mean-recovery
constraint is 4-momentum conservation: $\mathbf{v}_\text{phys} =
N^{-1}\sum V_{ijklm'n'} \in [0,1]^4$.
The physical 4-momentum is the invariant; each detector subsystem is a
frequency-band decomposition of the same event.

\subsection{Molecular spectral holography}

The spectral hologram $\Hologram(\omega,t) = \sum_n c_n(t)S_n(\omega)
e^{i\phi_n(t)}$ \citep{companion:hologram} superimposes ground, thermal,
and excited state spectra.  Mean-recovery: the three spectra are three
frequency-band decompositions of the same molecular oscillation; their
superposition recovers six quantities inaccessible to single-state measurement.

\subsection{Atomic spectroscopy (H, H$_2$, H$_2$O)}

The categorical spectrometer \citep{companion:hyperfine} uses four hardware
oscillators (CPU clock, memory bus, LED, refresh polarity) as four
frequency-band decompositions of the same atomic/molecular oscillation.
The 280/280 cross-validation (100\% agreement) is the empirical verification
of Theorem~\ref{thm:unified_signal}: all four decompositions of the same
physical state agree.

\subsection{Single-beam positioning}

A single electromagnetic signal simultaneously serves as radio-band,
microwave-band, and optical-band positioning reference \citep{companion:signal_source}.
Mean-recovery forces all three band-specific position estimates to converge
to the same physical position.  Theorem~\ref{thm:coherence} explains why:
the spectral separation between bands makes inconsistency exponentially
suppressed.

% =========================================================
\section{The Missing Constant}
\label{sec:constant}
% =========================================================

The Unified Signal Hypothesis paper conjectured: "the fundamental constant
of physics is not the speed of light but the universal mean-recovery
constraint."

\begin{theorem}[Mean-Recovery Is the Unique Universal Constraint]
\label{thm:fundamental}
The mean-recovery constraint $N^{-1}\sum_i \Sv_i = \Sv$ is the unique
constraint that:
(i) holds for all bounded oscillatory systems regardless of energy scale;
(ii) reduces to all known conservation laws in their respective domains;
(iii) is derivable from a single axiom (Bounded Phase Space);
(iv) has no free parameters.
\end{theorem}

\begin{proof}
(i) By Theorem~\ref{thm:poincare_mean_recovery}, mean-recovery is equivalent
to Poincar\'e recurrence, which holds for every bounded measure-preserving
system.  No energy-scale restriction applies.

(ii) In each domain:
MS/MS: mean-recovery = fragment mass conservation.
LHC: mean-recovery = 4-momentum conservation.
Holography: mean-recovery = energy conservation across electronic states.
Positioning: mean-recovery = uniqueness of physical position.

(iii) The Bounded Phase Space Axiom + Poincar\'e recurrence derives
mean-recovery with no additional postulates.

(iv) The constraint $N^{-1}\sum \Sv_i = \Sv$ has no adjustable
parameters: the mean is exactly the physical state.
\end{proof}

\begin{corollary}[Speed of Light as Derived]
The speed of light $c = 2\pi\,\text{rad/tick}$ (categorical angular units,
\citealt{companion:math}) is a projection of the partition traversal rate
onto the spatial measurement axis.  It is derived from the mean-recovery
constraint and the partition structure of three-dimensional bounded space;
it is not a primitive constant of nature.
\end{corollary}

% =========================================================
\section{Applications}
\label{sec:applications}
% =========================================================

\subsection{Frequency-agnostic trigger design}

The LHC trigger currently operates at one frequency (the 40\,MHz bunch
clock).  By the Unified Signal Theorem, the collision event simultaneously
encodes information at all four optical channel frequencies (10\,keV to MeV).
A trigger that combines all four channels is not using four independent
measurements --- it is using four frequency-band decompositions of the same
event.  The SEBD reconstruction of Section~7 of \citet{companion:extensor}
implements this exactly.

\subsection{Multi-domain MS instrument design}

A mass spectrometer that simultaneously reads out all four analyzer
frequencies (TOF, quadrupole, Orbitrap, FT-ICR) is not measuring four
things --- it is measuring the same partition state in four frequency-band
decompositions.  The stacked virtual partition tensor
$V_{ijkl}$ \citep{companion:tensor} is the instrument-level implementation
of the Unified Signal Theorem.

\subsection{Single-beam positioning}

The multi-domain positioning result of \citet{companion:signal_source}
(GPS-agnostic positioning from one cellular signal) is a direct application
of Theorem~\ref{thm:unified_signal}.  The theoretical foundation is now
complete: the cellular signal is a bounded oscillatory signal; its radio,
microwave, and optical interpretations are three frequency-band decompositions;
their agreement on position follows from mean-recovery.

\section*{Acknowledgements}
This work was supported by the AIMe Registry for Artificial Intelligence.

\bibliographystyle{plainnat}
\bibliography{references}

\end{document}
